\chapter{ Introdução}
\section{ Contextualização}
Na sua definição mais básica, uma base de dados é um conjunto de informações inter-relacionadas, utilizadas para armazenar e organizar dados de forma a facilitar a sua gestão e o acesso. Os dados referem-se a qualquer informação capturada sobre uma entidade e os seus atributos. À medida que uma coleção de dados cresce e assume maior complexidade, torna-se mais difícil manter a organização e a acessibilidade, surgindo a base de dados como a solução central para facilitar a interpretação desta informação.

Neste âmbito, os Sistemas de Base de Dados permitem-nos armazenar, recuperar e relacionar dados de forma estruturada. A plataforma MySQL, um dos sistemas mais utilizados, serve como interface de gestão para bases de dados relacionais através da linguagem SQL. No setor da saúde, estas ferramentas são fulcrais para garantir o funcionamento eficiente das instituições, nomeadamente na gestão de dispositivos médicos. A monitorização e manutenção destes equipamentos, apoiadas por bases de dados, são essenciais para o controlo operacional diário, garantindo que a tecnologia médica é utilizada com a máxima eficiência.

A manutenção regular é fundamental para assegurar a segurança dos pacientes e a qualidade dos serviços de saúde. Equipamentos em condições subótimas podem levar a diagnósticos incorretos ou tratamentos inadequados, colocando vidas em risco. Dada a complexidade dos aparelhos hospitalares, é imprescindível um sistema organizado que facilite o agendamento de manutenções preventivas e corretivas. Um sistema bem estruturado permite antecipar falhas, minimizar interrupções e garantir que os aparelhos cumpram as normas de segurança. Além disso, a rastreabilidade dos históricos permite identificar tendências de desgaste, otimizando recursos e reduzindo custos com reparos emergenciais ou substituições prematuras, promovendo, em última análise, a sustentabilidade operacional e financeira do hospital.

\newpage
\section{ Apresentação do Caso de Estudo}
O presente caso de estudo foca-se no desenvolvimento de um sistema para a gestão e acompanhamento da manutenção de equipamentos hospitalares. O sistema centraliza informações sobre os dispositivos, a sua localização geográfica nas unidades de saúde e os departamentos responsáveis.

A solução proposta integra os seguintes tópicos:

\begin{itemize}
    \item Gestão de Intervenções: Registo de ordens de serviço, históricos de manutenção e acompanhamento técnico.
\end{itemize}

\begin{itemize}
    \item Gestão de Recursos: Monitorização dos técnicos envolvidos, inventário de peças utilizadas e relação com fornecedores.
\end{itemize}

\begin{itemize}
    \item Controlo Financeiro: Rastreabilidade dos custos associados a cada intervenção e análise de desgaste dos ativos ao longo do tempo.
\end{itemize}

\section{ Motivação e Objetivos}
A motivação deste trabalho reside na aplicação prática dos conceitos de planeamento e desenvolvimento de bases de dados relacionais a um cenário real e complexo. O objetivo central é projetar uma solução que garanta a integridade e a acessibilidade dos dados hospitalares.

Para alcançar este propósito, o projeto divide-se nos seguintes objetivos específicos:

\begin{itemize}
    \item Modelação: Elaboração dos modelos conceptual (Diagrama Entidade-Associação), lógico e físico.
\end{itemize}

\begin{itemize}
    \item Implementação: Criação da base de dados em MySQL e respetivo povoamento com dados realistas.
\end{itemize}

\begin{itemize}
    \item Programação SQL: Desenvolvimento de subprogramas avançados, incluindo procedimentos (procedures), funções e triggers, para automatizar tarefas e facilitar a extração de métricas relevantes para a gestão hospitalar.
\end{itemize}